\section{Experiments}
\label{sec:experiments}
We report fixed-frame census summaries: medians over all non-leading targets (the mean's unique leader is excluded and costs zero under it), with the interquartile range in brackets, and, whenever a cap is binding, the feasible fraction next to every cost. Raw $K$ and $L_1$ are not comparable between the 6-task and 57-task frames. Medians and quartiles over targets use the inverted-CDF definition (for an even number of targets, such as the 78 HELM targets, the lower of the two middle values); the usual midpoint median is given in the text where it differs (for example 11.5 against 11 for the minimum rule's $K$ on HELM) and for every cell in Appendix Tables~\ref{tab:rangesceil}--\ref{tab:medconv} (with resampling ranges for the medians and feasible counts). A range in brackets after a count or share is the 2.5--97.5 percentile range over 1000 resamplings of the targets of the fixed census (seed 23): a sensitivity range for that census, not a confidence interval for a population of models; Wilson intervals, a binomial reference, are in the same file. Costs for every target, cap and rule are in \texttt{analysis/results/*.csv}; the table and figure to script map is in \texttt{figtab/README.md}.

\subsection{How much is needed, and how it varies with the rule}
\begin{table}[t]
\caption{Minimum investment to become uniquely first, all non-leading targets of each fixed census: median [Q1, Q3] (inverted-CDF convention: lower middle value for even $N$) of tasks touched $K$ and total score points $L_1$ under the ceiling cap (scores may rise to 100) and under a 10-point per-task cap. feas.: targets that can reach first place; $^\dagger$: median over feasible targets only (a different target set for each rule). Condorcet winner is a requirement, an upper bound for Condorcet-consistent rules. Ordinal rules: HELM and the top-500 frame only (separate block; the frames are not pooled). Source: \texttt{analysis/results/summary/main\_descriptive.csv}.}
\label{tab:main}
\centering\scriptsize\setlength{\tabcolsep}{3pt}
\input{../tables/tab_main_ceil_gain10_all_nonleader.tex}
\end{table}
\paragraph{Unrestricted headroom.}
With every score allowed to rise to 100 (Table~\ref{tab:main}, left), a typical non-leading Open LLM Leaderboard model needs to raise 2 of its 6 tasks and 181 points in total under the mean; the median is 3 tasks and 106 points, the minimum 3 tasks and 43 points. On HELM the mean needs 16 of 57 subjects and 661 points (midpoint median 696), the minimum 11 subjects (11.5) and 98 points, plurality 9 subjects and 47 points. Per task, the mean's typical HELM target needs 11.6 points on 28\% of the subjects, against 30.1 points on 33\% of the tasks on the Open LLM Leaderboard; raw units are not commensurable between the frames. Because the minimum rule only requires lifting a model's weakest tasks to the rival's weakest score, whereas the mean requires closing $n$ times the gap in the average, the \emph{points} needed are lower under every other rule we compare than under the mean for at least 98.7\% of targets (paired on the same target, Appendix Fig.~\ref{fig:paired}), on both frames, in raw points; in units of each task's standard deviation the mean no longer has the largest median $L_1$ among the score rules (geometric mean above it in all three frames, Appendix Table~\ref{tab:sdunits}). The number of \emph{tasks} tells a different story and depends on the frame: on HELM the minimum needs fewer tasks than the mean for 91\% of targets and plurality for 81\%, whereas the median needs more for 82\%, Borda for all targets and the Condorcet requirement for 87\%; on the Open LLM Leaderboard the minimum needs more tasks than the mean for 81\% of targets (equal for 19\%) and the median needs the same number for 53\%. These comparisons mix the rule's response to increases with the different gap each rule induces; the constructed examples (Table~\ref{tab:counter}) show that the sign can flip, so none of this is a universal ordering. Cost tends to grow with the target's rank for every rule (Appendix Fig.~\ref{fig:rank}).

\paragraph{The headroom assumption decides whether first place is reachable.}
With at most $+10$ points per task, most targets cannot reach first place under any rule (Table~\ref{tab:main}, right; Appendix Fig.~\ref{fig:caps}a): 113 of 4{,}495 Open LLM Leaderboard targets (2.5\% [2.0, 3.0]) under the mean, 446 (9.9\% [9.0, 10.8]) under the minimum; on HELM 29 of 78 (37\% [26, 49]) under the mean, 14 (18\% [10, 27]) under the minimum, but 57 (73\% [63, 83]) under plurality. The $+10$ cap is a scenario with no data-based justification; at $+5$ and $+20$ the mean reaches first place for 13 (0.3\%) and 830 (18\%) Open LLM Leaderboard targets and for 11 (14\%) and 56 (72\%) HELM targets (Appendix Table~\ref{tab:gain520}), so the feasibility percentages are properties of the cap and not of the leaderboards. Resampling the targets of each census with replacement (1000 draws, seed 23; analysis/results/paired\_target\_bootstrap.csv; a sensitivity range for the selected census, not a population confidence interval; Appendix Fig.~\ref{fig:bootpaired}) gives for the difference in feasible fraction at the 10-point cap, minimum minus mean, $-0.19$ [$-0.29$, $-0.09$] on HELM, $+0.07$ [$0.07$, $0.08$] on the 4,496-model Open LLM Leaderboard and $+0.47$ [$0.41$, $0.52$] on the top-500 frame;. At $+10$ the minimum is thus the most reachable score rule on the Open LLM Leaderboard and the least reachable on HELM (Section~\ref{sec:reversal} asks how far that is a property of the frames). By the threshold equivalence of Section~\ref{sec:setting}, the feasible fraction under a gain-$h$ cap is exactly the fraction of targets whose ceiling-cap $L_\infty$ is at most $h$, which is why $L_\infty$ is reported as a feasibility diagnostic. The empirical gain scales (\texttt{emp95/99}) are descriptive only.

\subsection{Does targeting matter?}
\label{sec:targeting}
A factor of 13--15 in tasks for the minimum rule against random order on 20 near-top HELM targets concerns a naive baseline on few targets, so we compare the optimum with greedy orders over \emph{all} feasible targets of each frame (analysis/results/final\_controls\_summary.csv; 110{,}876 target-rule-cap records; ranges from 1000 target resamplings shared by all policies). Raising tasks in order of lowest score first attains the optimal $K$ for every target at the ceiling cap (Figure~\ref{fig:controlsK}), for all five score rules and all three frames (15 panels), and at both caps it is optimal for the minimum rule, as highest-headroom-first is for the mean (Propositions~\ref{prop:mean}--\ref{prop:min}). Targeting thus beats random order, whose median task count over the 78 HELM targets at the ceiling is 1.76 [1.67, 1.92] (mean), 1.14 [1.08, 1.21] (median) and 4.23 [3.44, 6.00] (minimum) times the optimum, but shows no advantage in $K$ over a greedy rule that needs no optimiser. Under the 10-point cap the median rule is where the optimiser matters: lowest-score-first needs 2.6 [2.4, 2.8] times the optimal $K$ on HELM (1.7 on the Open LLM Leaderboard) and matches it for none of 40 HELM targets. In points, cap-prefix greedy policies have median $L_1$ 7.7 times (median rule) and 4.8 times (minimum) the optimum on HELM at the ceiling, against 1.0 under the mean (Appendix Fig.~\ref{fig:controlsL1}). These describe the named policies, not training costs.
\begin{figure}[t]
\centering\includegraphics[width=\linewidth]{../figures/fig11m_policy_controls_main.pdf}
\caption{Tasks needed by four policies relative to the optimal $K$ (log axis; median over all feasible targets, bars: 2.5--97.5 percentile range over 1000 target resamplings shared across policies; numbers of feasible targets are in Appendix Table~\ref{tab:controlsn}). Lowest-score-first coincides with the optimum at the ceiling for every score rule and frame. Source: \texttt{figures/fig11m\_policy\_controls\_main.provenance.json}.}
\label{fig:controlsK}
\end{figure}

\begin{table}[tb]
\caption{Which comparisons of the minimum rule (min) with the mean survive the checks. S: stable; R: reversed; ---: not changed by the check. Full frames: HELM, Open LLM Leaderboard (OLLB), top-500 (T500). Matched frames: random 79-model OLLB subsets, mean-top-79 frame, random six-subject HELM subsets (HELM-6); median or majority sign. SD units: costs per task standard deviation. IQR cap: cap at each task's interquartile range. Records which orderings change, not why. Sources: Table~\ref{tab:main}; Appendix Tables~\ref{tab:frames}, \ref{tab:capsmain}, \ref{tab:sdunits}, \ref{tab:framegroupsceil}, \ref{tab:framegroupsgain10}.}
\label{tab:verdict}
\centering\scriptsize\setlength{\tabcolsep}{3pt}
\begin{tabular}{lllll}
\toprule
comparison & HELM / OLLB / T500 & random OLLB-79 / top-79 / HELM-6 & SD units & IQR cap (HELM / OLLB / T500) \\
\midrule
Fewer tasks $K$ at ceiling & min / mean / mean & S / S / R & --- & --- \\
Lower $L_1$ at ceiling & min / min / min & S / S / S & S & --- \\
Mean costliest score rule in $L_1$ & yes / yes / yes & S / S / S & R & --- \\
More reachable at $+10$ & mean / min / min & S (84) / R / S (89) & --- & S / R / S \\
\bottomrule
\end{tabular}
\end{table}

\subsection{What the reversal depends on}
\label{sec:reversal}
At the 10-point cap the minimum rule is more reachable than the mean on the Open LLM Leaderboard and less reachable on HELM. Matching one dimension and the cap and unit change this picture (Figure~\ref{fig:frames}; 100 random frames per row, seed 31; ranges across selected frames are sensitivity ranges, not confidence intervals). Among 100 random 79-model subsets of the Open LLM Leaderboard (the number of rivals matched to HELM) the minimum is more reachable than the mean in 84 draws, less in 12 and equal in 4 (median paired difference $+0.08$ [$-0.05$, $+0.31$]); in the fixed top-79 frame selected by the mean it is less reachable (59 of 78 against 78 of 78, $-0.24$). Among 100 random six-subject subsets of HELM (the task count matched to the Open LLM Leaderboard) the minimum is less reachable in 89 draws and more in 9 ($-0.18$ [$-0.37$, $+0.23$]). So each sign survives matching in most random draws, with exceptions, and the mean-selected frame reverses it; random subsets keep the task correlation of their parent frame, mean-selected frames have almost none ($0.02$ and $0.08$; Appendix Table~\ref{tab:frames}), but the frames differ in several respects at once and none of this identifies a cause. The cap matters as much (Appendix Fig.~\ref{fig:capsfig}, Tables~\ref{tab:capsmain}, \ref{tab:capscounts}): under a cap equal to each task's interquartile range of scores (a descriptive scale, not an achievable gain) the full Open LLM Leaderboard reverses to the mean 582 against the minimum 283 of 4{,}495 targets, whereas the minimum leads on the top-500 frame (129 against 52 of 499) and the mean on HELM (39 against 20 of 78). The reachability ranking of the rules therefore depends on frame, cap and unit, and the $+10$ reversal is not a property of the leaderboards (Table~\ref{tab:verdict} lists which comparisons survive).


\begin{figure}[t]
\centering\includegraphics[width=\linewidth]{../figures/fig13_frames.pdf}
\caption{Minimum minus mean in the fraction of targets reachable at the 10-point cap, over frames matched in rivals (random 79-model subsets of the Open LLM Leaderboard, and the fixed mean-top-79 frame) or tasks (random six-subject subsets of HELM): (a) paired difference per frame, black: median and 2.5--97.5\% range over 100 frames (seed 31); counts: frames where the minimum is more, less or equally reachable; (b) the fractions; open: full frames. Ranges over selected frames are not confidence intervals. Source: \texttt{figures/fig13\_frames.provenance.json}; Appendix Tables~\ref{tab:frames}, \ref{tab:framegroupsgain10}.}
\label{fig:frames}
\end{figure}

\subsection{Sensitivity}
Resampling the tasks (100 draws, ten near-top targets per frame, cap $+10$) changes comparisons substantially: on HELM, whether the minimum rule can reach first place varies strongly across draws (feasible in 0.11 to 1.0 of the draws depending on the target, against at least 0.88 for every other rule). Score noise and the HELM task-count experiment give percentile ranges under stated assumptions, not confidence intervals (Appendix~\ref{app:sens}). On 20 near-top targets per frame, the margin $\varepsilon$ and the rival set also change task counts (Appendix Tables~\ref{tab:eps}, \ref{tab:frame}).

\paragraph{Ordinal rules on the top-500 Open LLM Leaderboard frame.} The ordinal rules are computed on a fixed top-500 frame (selected once by the normalised mean; 499 non-leading targets), kept separate from the 4,496-model census and from HELM (Table~\ref{tab:main}, separate block; Appendix Table~\ref{tab:top500}). At the ceiling the median target needs $K=4$ (Borda), 2 (plurality) and 1 (mean) tasks; at the 10-point cap the reachable targets are Borda 325, plurality 53, mean 113, minimum 348 and the Condorcet requirement 60 of 499, a ranking unlike HELM's.

\paragraph{Processing choices in the Open LLM Leaderboard frame.} The main frame keeps the largest reported average for each of 79 repeated names (74 of the 79 have a lower score on some task than the lowest-average duplicate) and drops one flagged model. Keeping the lowest average, including the flagged model, or facing only non-merged, unflagged rivals leaves the 10-point feasibility ordering of the score rules and the ceiling-cap $L_1$ order unchanged (reachable for the mean: 113 of 4{,}495 in the main frame, 111, 113, and 81 of 3{,}788 without merged models; Appendix Table~\ref{tab:processing}). The \emph{raw} score scale changes the 10-point ordering (reachable: harmonic 459, geometric 382, mean 296, minimum 224, median 21), as $+10$ raw points is a different intervention from $+10$ normalised points. The mean-selected top-500 frame lowers the mean's cost by construction ($L_1=71$ against $181$).
